%% ============================================================================
%%  CV – Georges Parfait Djimefo Kapen (VERSION FRANÇAISE)
%%  Adapté pour : Développeur IA (Edmonton) — Deloitte Omnia AI Strategy
%%  Dernière mise à jour : Avril 2026
%% ============================================================================
\documentclass[10pt,a4paper,sans]{moderncv}

\moderncvstyle{banking}
\moderncvcolor{blue}

\usepackage[scale=0.88, top=1.0cm, bottom=1.0cm]{geometry}
\setlength{\hintscolumnwidth}{2.8cm}

\usepackage[utf8]{inputenc}
\usepackage[T1]{fontenc}
\usepackage{lmodern}
\usepackage{enumitem}
\usepackage{hyperref}
\hypersetup{colorlinks=true, urlcolor=blue!70!black}

% ============================================================================
%  DONNÉES PERSONNELLES
% ============================================================================
\name{Georges Parfait}{Djimefo Kapen}
\title{Développeur IA \textbar{} IA Agentique \& Systèmes Infonuagiques}
\address{Montréal, Québec, Canada}{}{}
\phone[mobile]{+1~514-245-3182}
\email{gpdk2025@gmail.com}

% ============================================================================
\begin{document}
\makecvtitle
\vspace{-1.8em}

% ============================================================================
%  PROFIL
% ============================================================================
\section{Profil}
Développeur IA avec expérience pratique en conception et déploiement de
\textbf{systèmes d'IA agentiques en production}, \textbf{flux de travail
multi-agents basés sur des LLM} et \textbf{agents conversationnels RAG
à l'échelle de l'entreprise} sur AWS.
Solide expertise en \textbf{évaluation et analyse comparative de modèles d'IA},
\textbf{architecture infonuagique native} (sans serveur, conteneurs),
\textbf{automatisation CI/CD} et \textbf{pipelines de données ETL/ELT}.
Doctorat prévu été~2026 avec recherche en \textbf{IA},
\textbf{réseaux de neurones sur graphes}, \textbf{apprentissage par renforcement
multi-agents} et \textbf{cybersécurité}.
Aptitude démontrée à communiquer les compromis techniques aux cadres dirigeants
et aux parties prenantes non techniques. Innovation, travail d'équipe,
leadership et orientation vers les résultats.

% ============================================================================
%  COMPÉTENCES TECHNIQUES
% ============================================================================
\section{Compétences techniques}
\cvitem{LLM\,/\,Agents}{LangChain, LangGraph, LangSmith, MCP, RAG (FAISS), CRAG, RAGAS, Claude (Haiku, Opus), Kiro}
\cvitem{IA\,/\,ML}{PyTorch, TensorFlow, Scikit-Learn, MARL (QMIX, QTRAN), GNN, DQN, NLP, Vision par ordinateur, LLM}
\cvitem{Infonuagique}{AWS (SageMaker, Bedrock, EKS, Lambda, CloudWatch, S3, ECR, CloudFormation), OpenStack, Kubernetes, Docker}
\cvitem{DevOps\,/\,IaC}{CI/CD (GitHub Actions, CodeBuild, OIDC), Git, Jenkins, GitLab, Terraform, Infrastructure en tant que code}
\cvitem{Données}{PySpark, Databricks, Snowflake, Delta Lake, SQL, Bases de données relationnelles, ETL/ELT}
\cvitem{Langages}{Python, Java, C++, PHP, SQL, R, JavaScript, TypeScript}
\cvitem{Web\,/\,API}{REST APIs, FastAPI, Streamlit, React, Angular, Node.js, D3.js, Plotly, Dash, Power~BI, Tableau}
\cvitem{Méthodes}{Agile/Scrum, TDD, Patrons de conception, Modélisation statistique, MLOps, Gouvernance IA}

% ============================================================================
%  EXPÉRIENCE PERTINENTE
% ============================================================================
\section{Expérience pertinente}

\cventry{2026}{Stagiaire ML --- IA Agentique \& Systèmes Infonuagiques}{Ericsson}{Montréal}{}{%
\begin{itemize}[leftmargin=1.2em, topsep=1pt, itemsep=0pt, parsep=0pt]
  \item Conçu et déployé un \textbf{système autonome d'IA agentique RAG en production}
        qui automatise la détection d'anomalies dans les journaux d'unités
        radio en télécommunications, remplaçant la revue manuelle pour les
        clients entreprise.
  \item Construit un \textbf{flux de travail multi-agents LLM} : agent ReAct à
        6~outils orchestré par Claude Haiku~4.5 sur Bedrock (notation
        SageMaker, RAG FAISS, modèle physique RF, classification de
        motifs, surveillance CloudWatch), reproduisant le raisonnement
        d'un analyste.
  \item Évalué et comparé les modèles fondamentaux avec le score de
        fidélité/pertinence \textbf{RAGAS} et une notation de confiance
        \textbf{méta-cognitive} fondée sur des preuves (ÉLEVÉ/MOYEN/FAIBLE)
        pour optimiser les compromis performance-coût.
  \item Entraîné un \textbf{SLM à double attention} personnalisé atteignant un
        \textbf{F1 de 90,6\,\%}, une précision de 99,97\,\% et une latence
        d'inférence de \textbf{1\,ms/journal} sur AWS SageMaker.
  \item Conçu un \textbf{pipeline map-reduce à l'échelle de la production} :
        triage $\rightarrow$ regroupement par signature de panne $\rightarrow$
        investigation parallèle par agents $\rightarrow$ synthèse adaptative
        inter-grappes. Traite 120+ anomalies en $\sim$12\,min.
  \item Construit un \textbf{tableau de bord Streamlit pour les dirigeants}
        exposant les métriques de précision, d'utilisation, de latence et
        de coût par requête pour les décisions de la direction.
  \item Déployé en tant que \textbf{pod K8s sur EKS} via tunnel Cloudflare ;
        construit un \textbf{pipeline événementiel} (S3 $\rightarrow$ Lambda
        $\rightarrow$ SQS $\rightarrow$ worker K8s $\rightarrow$ S3
        $\rightarrow$ alertes SNS) pour la haute disponibilité.
  \item Conçu une \textbf{boucle de réentraînement entièrement autonome} :
        déclenchée par la dérive + hebdomadaire via EventBridge, avec
        promotion automatique uniquement lorsque le F1 du nouveau modèle
        dépasse la référence en production (gouvernance et surveillance IA).
  \item Mis en \oe{}uvre un \textbf{pipeline CI/CD en 7~étapes} (GitHub Actions +
        CodeBuild avec OIDC) : test $\rightarrow$ analyse de sécurité OWASP
        $\rightarrow$ infra CloudFormation $\rightarrow$ validation pré-prod
        $\rightarrow$ déploiement $\rightarrow$ test de fumée $\rightarrow$
        retour arrière automatique. Zéro clé IAM statique.
  \item Rédigé \textbf{131 tests automatisés} avec un gardien de tests adaptatif
        basé sur l'AST assurant une couverture de 100\,\% des fonctions
        en production.
  \item \textit{Pile technologique :} Python, AWS (SageMaker, Bedrock, EKS,
        Lambda, SQS, SNS, S3, ECR, CloudFormation, CloudWatch), LangChain,
        LangGraph, LangSmith, FAISS, MCP, Claude, Kiro, Docker, Kubernetes,
        CI/CD, Terraform.
\end{itemize}}

\cventry{2022--2024}{Développeur IA~I --- Systèmes de données d'entreprise}{Intact Corporation financière}{Montréal}{}{%
\begin{itemize}[leftmargin=1.2em, topsep=1pt, itemsep=0pt, parsep=0pt]
  \item Ingérer et traiter des données à grande échelle.
  \item Surveiller les pipelines et développer des solutions de tests automatisés
        assurant la fiabilité des systèmes.
  \item \textit{Pile technologique :} PySpark, SQL, Databricks, Snowflake, AWS,
        Shell, Git.
\end{itemize}}

% ============================================================================
%  EXPÉRIENCE ADDITIONNELLE
% ============================================================================
\section{Expérience additionnelle}

\cventry{2025--2026}{Chargé de cours}{Polytechnique Montréal}{Montréal}{}{%
\begin{itemize}[leftmargin=1.2em, topsep=1pt, itemsep=0pt, parsep=0pt]
  \item Enseigner le cours de programmation procédurale Python ; communiquer
        des concepts techniques complexes à des auditoires non techniques.
\end{itemize}}

\cventry{2022--2023}{Chargé de cours}{Polytechnique Montréal}{Montréal}{}{%
\begin{itemize}[leftmargin=1.2em, topsep=1pt, itemsep=0pt, parsep=0pt]
  \item Enseigner le cours de programmation procédurale Python.
\end{itemize}}

\cventry{2021--2022}{Auxiliaire scientifique --- MOOC IA}{Université de Montréal}{Montréal}{}{%
\begin{itemize}[leftmargin=1.2em, topsep=1pt, itemsep=0pt, parsep=0pt]
  \item Concevoir des cours et travaux pratiques en apprentissage supervisé
        (Python).
\end{itemize}}
\cventry{2012--2014}{Ing\'{e}nieur statisticien \'{e}conomiste contractuel}{Minist\`{e}re du Tourisme et des Loisirs}{Cameroun}{}{%
\begin{itemize}[leftmargin=1.2em, topsep=1pt, itemsep=0pt, parsep=0pt]
  \item Analyser des donn\'{e}es
\end{itemize}}

\cventry{2010--2012}{Professeur associ\'{e} de statistique et probabilit\'{e}}{Universit\'{e} Protestante d'Afrique Centrale}{Cameroun}{}{%
\begin{itemize}[leftmargin=1.2em, topsep=1pt, itemsep=0pt, parsep=0pt]
  \item Enseigner les cours de statistique et de probabilit\'{e}.
\end{itemize}}

\cventry{2009--2010}{Charg\'{e} de cours en analyse de donn\'{e}es}{ISSEA (Institut Sous-r\'{e}gional de Statistique et d'\'{E}conomie Appliqu\'{e}e)}{Cameroun}{}{%
\begin{itemize}[leftmargin=1.2em, topsep=1pt, itemsep=0pt, parsep=0pt]
  \item Enseigner le cours d'analyse de donn\'{e}es (R).
\end{itemize}}
% ============================================================================
%  PUBLICATIONS
% ============================================================================
\section{Publications}

\textbf{[1]} \textit{Energy-Efficient and Near Real-Time Routing Algorithm for
Disaster Monitoring in Wireless Sensor Networks}\\
G.\,P.~Djimefo~Kapen, R.~Al~Mallah, S.~Pierre --- Polytechnique
Montréal \hfill \textit{(Publié)}
\begin{itemize}[leftmargin=1.5em, topsep=2pt, itemsep=0pt, parsep=0pt]
  \item Conçu \textbf{GAPF}, un cadre novateur de \textbf{méta-apprentissage}
        fusionnant des politiques MARL (QMIX \& QTRAN) via un
        \textbf{encodeur GCN} avec seulement \textbf{1\,473 paramètres}.
  \item Atteint $\geq$98\,\% de livraison de paquets, 2--3,5$\times$ moins
        d'énergie, 13$\times$ moins de mémoire que les référentiels.
\end{itemize}

\vspace{0.4em}
\textbf{[2]} \textit{Carbon-Aware Autonomous Data Reduction and Self-Healing
Routing for 6G-Integrated Disaster Sensor Networks}\\
G.\,P.~Djimefo~Kapen, R.~Al~Mallah, S.~Pierre --- Polytechnique
Montréal \hfill \textit{(Soumis)}
\begin{itemize}[leftmargin=1.5em, topsep=2pt, itemsep=0pt, parsep=0pt]
  \item Proposé \textbf{CALASH}, le \textbf{premier protocole} optimisant
        conjointement l'intensité carbone du cycle de vie via un moteur
        hybride \textbf{Lyapunov--DQN}.
  \item Atteint 60\,\% de durée de vie réseau plus longue et 33\,\% de LCI
        inférieure aux référentiels.
\end{itemize}

\vspace{0.4em}
\textbf{[3]} \textit{Zero-Trust Shielded Graph-Structured Decision Control for
Secure Autonomous Recovery in AI-Native 6G Networks}\\
G.\,P.~Djimefo~Kapen, R.~Al~Mallah, S.~Pierre --- Polytechnique
Montréal \hfill \textit{(Soumis)}
\begin{itemize}[leftmargin=1.5em, topsep=2pt, itemsep=0pt, parsep=0pt]
  \item Conçu \textbf{CASTER-ZT}, combinant la prise de décision basée sur
        les \textbf{GNN} avec un bouclier d'admissibilité \textbf{zéro
        confiance} formellement borné ; sécurité et conformité d'entreprise
        par conception.
  \item Atteint 74--98\,\% de détection de nœuds malveillants avec
        \textbf{zéro faux positif}, latence de 3\,ms dans le budget
        O-RAN de 10\,ms.
\end{itemize}

% ============================================================================
%  PROJETS SÉLECTIONNÉS
% ============================================================================
\section{Projets sélectionnés}

\cventry{2022}{Stagiaire en ingénierie des données}{Intact Corporation financière}{Montréal}{}{%
\begin{itemize}[leftmargin=1.2em, topsep=1pt, itemsep=0pt, parsep=0pt]
  \item Manipuler des données (PySpark, SQL, Databricks, Snowflake, AWS).
\end{itemize}}

\cventry{2021}{Projet intégrateur --- IATA}{Polytechnique Montréal}{Montréal}{}{%
\begin{itemize}[leftmargin=1.2em, topsep=1pt, itemsep=0pt, parsep=0pt]
  \item « Flight Data eXchange --- Développer un outil initial et contextuel
        d'évaluation des données de vol. » (Dash, Python, React, Electron,
        Axios).
\end{itemize}}

\cventry{2020--2021}{Stagiaire en initiation à la recherche}{Polytechnique Montréal}{Montréal}{}{%
\begin{itemize}[leftmargin=1.2em, topsep=1pt, itemsep=0pt, parsep=0pt]
  \item Modélisation de reconnaissance vocale pour paiements bancaires
        sécurisés ; apprentissage profond (Python) déployé sur Android (Java).
\end{itemize}}

\cventry{2021}{Projet de TAL}{Polytechnique Montréal}{Montréal}{}{%
\begin{itemize}[leftmargin=1.2em, topsep=1pt, itemsep=0pt, parsep=0pt]
  \item Extraction de mots-clés à partir de textes ainsi que leur type
        (Python).
\end{itemize}}

\cventry{2021}{Projet de visualisation de données --- Chaire de recherche en fiscalité}{Polytechnique Montréal}{Montréal}{}{%
\begin{itemize}[leftmargin=1.2em, topsep=1pt, itemsep=0pt, parsep=0pt]
  \item Application web pour comparer les taxes et taux par année
        (Plotly.js, D3).
\end{itemize}}

\cventry{2020}{Entrepreneur / Stagiaire en initiation à la recherche}{Polytechnique Montréal}{Montréal}{}{%
\begin{itemize}[leftmargin=1.2em, topsep=1pt, itemsep=0pt, parsep=0pt]
  \item Modélisation de reconnaissance faciale intégrant les biais raciaux
        et de genre ; apprentissage profond (Python).
\end{itemize}}

\cventry{2019--2020}{Stagiaire en programmation}{FLEX GROUP}{Laval}{}{%
\begin{itemize}[leftmargin=1.2em, topsep=1pt, itemsep=0pt, parsep=0pt]
  \item Modélisation de reconnaissance optique de caractères automatique
        (Python, PHP).
\end{itemize}}

\cventry{2006}{Stagiaire pour maîtrise en statistique appliquée}{Banque Internationale du Cameroun pour l'Épargne et le Crédit}{Cameroun}{}{%
\begin{itemize}[leftmargin=1.2em, topsep=1pt, itemsep=0pt, parsep=0pt]
  \item Modèles mathématiques de prévision des recettes et dépenses
        d'une banque (R).
\end{itemize}}

% ============================================================================
%  FORMATION
% ============================================================================
\section{Formation}

\cventry{Été 2026}{Doctorat en génie informatique}{Polytechnique Montréal}{Montréal, Canada}{}{%
\begin{itemize}[leftmargin=1.2em, topsep=1pt, itemsep=0pt, parsep=0pt]
  \item Routage piloté par l'IA, réseautique écoresponsable et sécurité
        zéro confiance pour les réseaux de capteurs 6G.
\end{itemize}}

\cventry{}{Diplôme d'ingénieur en génie logiciel}{Polytechnique Montréal}{Montréal, Canada}{}{%
\begin{itemize}[leftmargin=1.2em, topsep=1pt, itemsep=0pt, parsep=0pt]
  \item Concentration : Intelligence artificielle --- Science des données.
\end{itemize}}

\cventry{}{Maîtrise en statistique appliquée}{École Nationale Supérieure Polytechnique de Yaoundé (ENSPY)}{Yaoundé, Cameroun}{}{%
\begin{itemize}[leftmargin=1.2em, topsep=1pt, itemsep=0pt, parsep=0pt]
  \item Équivalence : Maîtrise en probabilités --- statistiques.
\end{itemize}}

\cventry{}{Maîtrise en mathématiques}{Université de Yaoundé~I}{Yaoundé, Cameroun}{}{%
\begin{itemize}[leftmargin=1.2em, topsep=1pt, itemsep=0pt, parsep=0pt]
  \item Équivalence : Diplôme d'études supérieures spécialisées (DESS)
        en mathématiques.
\end{itemize}}

% ============================================================================
%  CERTIFICATIONS
% ============================================================================
\section{Licences et certifications}

\cvlistitem{AWS Certified Cloud Practitioner (CLF-C01)\newline
\textit{Parcours de préparation à la certification (Pluralsight), Montréal}}
\cvlistitem{AWS Certified Developer --- Associate (DVA-C02)\newline
\textit{Parcours de préparation à la certification (Pluralsight), Montréal}}

% ============================================================================
%  CONFÉRENCES
% ============================================================================
\section{Conférences et séminaires}

\cventry{Mars--Avr.\ 2026}{Présentation orale}{PolyCongrès 2026, Polytechnique Montréal}{Montréal, Canada}{}{« Energy-Efficient and Near Real-Time Routing for Disaster Monitoring in Wireless Sensor Networks --- Graph-Attentive Policy Fusion (GAPF). »}

\cventry{Oct.\ 2013}{Présentation orale}{20\textsuperscript{e} Conf.\ de l'Association africaine des scientifiques de l'insecte}{Yaoundé, Cameroun}{}{« Biogéographie des communautés d'insectes à l'aide de méthodes de regroupement. »}

\cventry{Août 2008}{Affiche}{29\textsuperscript{e} Conf.\ annuelle de la Société internationale de biostatistique clinique}{Copenhague, Danemark}{}{« Estimation de l'incidence du VIH à partir d'enquêtes transversales répétées. »}

\cventry{Jan.\ 2008}{Présentation orale}{ANMSA (Réseau Africain de Statistique Mathématique et ses Applications)}{Franceville, Gabon}{}{« Épidémiologie du VIH au Cameroun : estimation de l'incidence à partir des données de prévalence. »}

% ============================================================================
%  BÉNÉVOLAT
% ============================================================================
\section{Bénévolat}
\cvlistitem{Bénévole à la branche IEEE de Polytechnique Montréal, Montréal, 2018.}
\cvlistitem{Bénévole à la conférence ConFoo, Montréal, 2018.}

% ============================================================================
\end{document}
